\documentclass[preprint,12pt,authoryear]{elsarticle}

\usepackage{amsmath,amssymb}
\usepackage{booktabs}
\usepackage{graphicx}
\usepackage{threeparttable}
\usepackage{caption}
\usepackage{array}
\usepackage{float}
\usepackage{hyperref}
\setlength{\emergencystretch}{3em}

\journal{Journal of Asian Economics}

\begin{document}

\begin{frontmatter}

\title{RCEP and Supply Chain Relationships of Chinese Firms: Evidence from Relationship-Pair Data}

\author{Anonymous manuscript}

\begin{abstract}
We examine how the Regional Comprehensive Economic Partnership (RCEP) is associated with the international supply-chain relationships of Chinese listed firms. Using FactSet Revere disclosures, we construct a balanced panel of 630,280 relationship--year observations for 78,785 directed customer and supplier pairs involving 5,689 firms over 2017--2024. A difference-in-differences specification with firm--partner--direction pair fixed effects estimates a 1.36 percentage-point increase in the probability that an RCEP relationship is active at year-end relative to relationships with other foreign partners. The estimate is similar under member-specific implementation dates and several alternative samples. Partner-country--year evidence is consistent with three policy-linked pathways: tariff relief, a larger intermediate-goods share in China's imports from member economies, and a lower partner-country import share. The response is concentrated in supplier relationships, for which the estimated differential increase is 2.95 percentage points; the estimate for customer relationships is statistically indistinguishable from zero. The joint Wald test does not reject equality of the four pre-treatment coefficients to zero at the 5\% level ($p=0.078$), but does reject at the 10\% level. The results point to an economically meaningful relationship-level response while also indicating the limits of identification over a short post-entry period.
\end{abstract}

\begin{keyword}
RCEP \sep supply chains \sep relationship-pair panel \sep difference-in-differences \sep trade policy
\end{keyword}

\end{frontmatter}

\section{Introduction}
\label{sec:introduction}

Production in Asia is organized through dense cross-border networks in which firms repeatedly source inputs from, and sell output to, partners in neighboring economies. Recent evidence points to considerable persistence in the region's aggregate trade network even after a large common shock \citep{kiyota2022covid}, while new measures of regional value-chain position emphasize the importance of tracing production links beyond gross trade flows \citep{he2023regional}. The entry into force of the Regional Comprehensive Economic Partnership (RCEP) on January 1, 2022 began bringing 15 Asia-Pacific economies under a common trade framework. Its combination of coordinated tariff schedules, regional cumulation of origin, and trade-facilitation commitments can change both the cost of a cross-border transaction and the set of partners with which that transaction is feasible. Whether this institutional integration changes relationships between firms, rather than only aggregate trade, remains an open empirical question.

This paper studies that question using relationship-level data. We convert disclosed customer and supplier records in the FactSet Revere database into a balanced annual panel of 630,280 observations for 78,785 directed firm--partner pairs. The empirical design compares relationships with partners in RCEP member economies to relationships with other foreign partners before and after the common 2022 treatment date. The specification includes firm--partner--direction pair fixed effects and year fixed effects, with standard errors clustered by relationship pair.

The main estimate is 0.0136 (standard error 0.0040), implying a 1.36 percentage-point differential increase in the probability that a relationship is active at year-end. The estimate is 0.0144 (standard error 0.0041) when member-specific entry dates are used and remains positive in several alternative time windows and sample exclusions. The response is concentrated in supplier relationships, for which the estimated increase is 2.95 percentage points; the corresponding estimate for customer relationships is small and statistically insignificant. At the partner-country--year level, RCEP exposure is associated with tariff relief, a higher intermediate-goods share in China's imports from member economies, and a lower partner-country share of China's imports. These aggregate trade patterns, together with the supplier-side relationship response, are consistent with stronger regional production-network integration helping Chinese firms maintain or establish disclosed supplier relationships. The joint pre-treatment Wald test does not reject at the 5\% level but does reject at the 10\% level ($p=0.078$). We therefore interpret the estimates as a positive relative response while keeping the limitations of both the timing evidence and the comparison group explicit.

The paper is positioned against several closely related bodies of work. Research on trade agreements has established effects on bilateral trade, firm outcomes, and third-country adjustment \citep{baier2007free,sun2022rtaquality,lee2023thirdcountry}, while recent RCEP studies evaluate regional value chains in a GVC-CGE framework \citep{sun2023estimating} or examine how financing constraints relate to manufacturers' export partnerships with RCEP economies \citep{peng2025financing}. A second literature studies firms' use of preferences and the sourcing consequences of rules of origin \citep{hayakawa2023firm,tanaka2022rooexports,head2024laffer}. A third analyzes how trade-policy changes propagate through supplier search and firm-to-firm production links \citep{garg2023within,grossman2024tariffs}. These studies motivate a relationship-level test but do not establish whether RCEP's entry into force altered the annual activity of directed customer and supplier links.

Relative to this work, the paper makes three contributions. First, it shifts the outcome from aggregate trade or export participation to the annual activity of disclosed firm--partner relationships and estimates the policy response within the same directed pair over time. Second, it relates partner-country--year changes in tariff relief, the composition of intermediate-goods imports, and partner-country import shares to relationship activity, while recognizing that these aggregate measures do not observe firm-level preference claims, certificates of origin, customs transactions, or relationship values. Third, it separates supplier from customer links and examines how the initial relationship response varies with pre-policy relationship maturity, prior trade-agreement coverage, and manufacturing exposure.

The remainder of the paper is organized as follows. Section~\ref{sec:theory} reviews the literature and develops the hypotheses. Section~\ref{sec:design} describes the sample, variables, and empirical design. Section~\ref{sec:empirical} presents the baseline estimates, mechanism analysis, and robustness tests. Section~\ref{sec:heterogeneity} examines heterogeneous responses. Section~\ref{sec:further} provides further evidence on a competing explanation. Section~\ref{sec:conclusion} concludes with policy implications.

\section{Literature review and hypothesis development}
\label{sec:theory}

\subsection{Literature review}

The first strand of related research examines how trade agreements and trade institutions affect bilateral trade and firm behavior. Gravity-based evidence establishes that free-trade agreements increase trade between members \citep{baier2007free}. Recent JAE studies extend the analysis from trade flows to export-product quality and foreign direct investment, examine environmental provisions as a determinant of firm-level quality adjustment, and estimate employment responses to regional integration \citep{sun2022rtaquality,narayan2022agreements,zhu2025environmental,li2025integration}. The incidence of an agreement nevertheless depends on the institutional details faced by firms. \citet{hayakawa2023firm}, for example, measure the firm-level utilization of regional trade agreements from the importer's perspective, while \citet{ni2022standards} analyze heterogeneous trade effects across sectors. Other recent work examines how tariff costs and trade-policy uncertainty shape cross-border organization, export behavior, and input choice \citep{xie2023tariff,zhou2022trade,handley2024input,poilly2025regional}. Firm-level evidence on third-country effects further shows why adjustment outside an agreement must be part of the comparison \citep{lee2023thirdcountry}. This literature implies that membership and statutory tariff schedules alone are insufficient measures of firms' effective exposure.

The second strand studies regional production networks and RCEP more directly. Recent work constructs measures of regional dual value chains \citep{he2023regional}, evaluates RCEP through a GVC-CGE model \citep{sun2023estimating}, and examines the reorganization of export-product structures as global value chains change \citep{zhang2024restructuring}. Complementary JAE studies examine domestic value added in exports, GVC upgrading after foreign-capital withdrawal, financial and institutional determinants of value-chain participation, and risk-related changes in regional trade structure \citep{du2023higher,wu2023gvc,xu2024financial,ali2026governance,liu2026risks}. At the firm level, \citet{peng2025financing} study financing constraints and export partnerships with RCEP economies. Other evidence connects lower cross-border mobility costs to bilateral value-chain linkages \citep{ma2024mobility} and digital development to economies' positions in the intermediate-goods trade network \citep{huang2026intermediate}. Work on data restrictions and digital-trade provisions further illustrates that contemporary regional integration operates through non-tariff rules as well as customs duties \citep{ma2024data,li2026digitalrules}. These contributions establish the relevance of regional integration for production networks, but aggregate value-chain measures, simulations, and exporter-country links do not reveal whether a particular disclosed customer or supplier relationship remains active after implementation.

The third strand concerns rules of origin and the organization of input sourcing. Rules of origin determine whether firms can claim preferences and can therefore either promote regional sourcing or impose compliance costs. \citet{tanaka2022rooexports} find that relaxing origin requirements produced heterogeneous export responses across developing-economy beneficiaries. Studying reforms in Japan, \citet{hayakawa2023reform} document an increase in imports from least-developed economies but heterogeneous effects on domestic producers. Theoretically, \citet{head2024laffer} show that tighter content requirements initially expand regional parts sourcing but reduce it once the requirement passes a threshold. The design of these rules is itself endogenous to negotiation and differs between intermediate and final products \citep{hayakawa2024negotiating}; overlapping preference regimes add another layer to firms' compliance choices \citep{hayakawa2022multiple}. Recent JAE evidence also links rules of origin to imported inputs and firm innovation \citep{zhou2025origin}. Taken together, this work makes input sourcing central to the expected effect of RCEP's regional-cumulation provision.

A final set of studies highlights that supply chains are networks of durable relationships rather than anonymous spot transactions. Administrative shipment data reveal that ownership boundaries, distance, relationship specificity, and input characteristics shape where firms source products \citep{garg2023within}. Models of unexpected tariffs similarly predict renegotiation with incumbent suppliers or costly search for replacement partners \citep{grossman2024tariffs}. Empirical studies connect imported inputs and supply-chain links to firm productivity \citep{newman2023imports}, show that shocks can propagate through observed firm networks \citep{crosignani2022pirates}, and examine supply-chain spillovers in corporate innovation \citep{ou2026spillovers}. The present paper joins this relationship-based perspective to the RCEP literature. Its empirical unit preserves the firm, partner, and relationship direction, allowing the analysis to distinguish changes in relationship activity from changes in trade values or aggregate network position.

\subsection{Hypothesis development}

RCEP was signed on November 15, 2020 and entered into force for the first group of members on January 1, 2022. The agreement covers China, Japan, South Korea, Australia, New Zealand, and the ten ASEAN economies. It combines market-access commitments with rules on origin, customs procedures, services, investment, intellectual property, and electronic commerce \citep{shen2020rcep,ping2021rcep}. Recent research on digital-trade rules in regional agreements underscores the network implications of this broader institutional scope \citep{li2026digitalrules}. RCEP's common origin regime is especially relevant in Asia, where firms previously navigated several overlapping ASEAN-plus-one agreements whose product-level rules reflected different negotiating interests \citep{hayakawa2022multiple,hayakawa2024negotiating}.

The tariff schedule is phased. Approximately 65 percent of tariff lines receive immediate elimination, while other lines are reduced over five to twenty years and a small set is excluded. More than 90 percent of tariff lines are scheduled eventually to reach zero, but the timing differs by product and member \citep{liu2022rcep}. The 2022--2024 observation window therefore captures the initial institutional response and only the early portion of tariff liberalization; it should not be read as an estimate of the agreement's full long-run tariff effect.

The value of a supply-chain relationship reflects more than the current tariff on a shipment. Establishing a link requires partner search, screening, negotiation, and adaptation to relationship-specific production needs. Once formed, the link remains active when the expected surplus from continued trade exceeds the cost of renegotiation or replacement. This logic is consistent with evidence that relationship specificity and input characteristics shape firms' sourcing choices \citep{garg2023within} and with models in which tariff shocks induce renegotiation with incumbent suppliers or costly search for alternatives \citep{grossman2024tariffs}.

RCEP raises the expected surplus from qualifying intraregional trade through phased tariff reductions, common procedures, and regional cumulation. Value added from any member can count toward the originating status of a product, subject to the agreement's regional-value-content rules. The provision broadens the set of production arrangements capable of receiving preferential treatment and reduces the need to satisfy separate origin systems for each bilateral route. Utilization is not automatic---firms must still find compliance worthwhile \citep{hayakawa2023firm}---but the joint reduction in trade and coordination costs can both sustain an existing link and make a newly disclosed link worthwhile. This reasoning yields the first hypothesis:

\noindent\textbf{H1.} \textit{RCEP's entry into force increases the probability that a disclosed supply-chain relationship between a Chinese listed firm and a partner in an RCEP economy is active.}

The agreement can sustain relationships through three related channels. First, incremental tariff relief raises the net return to transactions with partners for which RCEP improves on the pre-existing tariff schedule. This margin should be strongest where the additional preference is economically meaningful, but its effect depends on actual utilization rather than the statutory rate alone. Second, regional cumulation increases the value of intermediate inputs obtained within the agreement because content contributed by one member can help a downstream product qualify in another. Evidence that origin rules affect export responses \citep{tanaka2022rooexports,hayakawa2023reform} and theoretical results linking content requirements to regional parts sourcing \citep{head2024laffer} support this input-reallocation channel. Recent firm-level work that connects origin rules with imported inputs provides a closely related motivation \citep{zhou2025origin}.

Third, a common origin regime can change the concentration of eligible sourcing. Under separate bilateral agreements, a firm may organize production around a narrow set of partners that satisfy route-specific requirements. Cumulation makes inputs from multiple member economies mutually usable for origin purposes, widening the feasible regional supplier set and reducing dependence on any one partner source. Research on import diversification under a regional cooperation initiative provides a related application of this sourcing margin \citep{zhu2025diversifying}. This prediction does not require every firm to diversify immediately: origin rules may generate nonlinear responses \citep{head2024laffer}, and reforms have produced heterogeneous effects across products and firms \citep{tanaka2022rooexports,hayakawa2023reform}. It instead implies that, on average, tariff relief, stronger intermediate-goods trade integration among members, and a lower partner-country import share are observable implications of the proposed channels. These trade measures are evaluated at the partner-country--year level. Their relationship-level implication is that deeper regional production integration can make supplier ties more valuable to maintain or establish. These arguments yield the second hypothesis:

\noindent\textbf{H2.} \textit{The relationship response to RCEP is associated with incremental tariff relief, stronger intermediate-goods trade integration among members, and reduced dependence on any one partner-country source.}

The relevance of these channels should differ by relationship direction. A supplier relationship records a foreign partner providing inputs or products to the Chinese firm, whereas a customer relationship records sales by the Chinese firm to a foreign partner. Tariff preferences can affect both directions, but regional cumulation operates first through the origin, eligibility, and organization of inputs. The supplier side is therefore exposed directly to the decision margins emphasized by research on imported-input choice \citep{handley2024input}, within-firm and external sourcing \citep{garg2023within}, and the productivity role of imports transmitted through supply chains \citep{newman2023imports}. Customer relationships may respond only after changes in input composition, production costs, or market access pass through to downstream sales. Over the first three years of implementation, this ordering implies a larger and more readily detectable response among supplier links. We therefore propose the third hypothesis:

\noindent\textbf{H3.} \textit{The increase in relationship activity following RCEP is stronger for supplier relationships than for customer relationships.}

\begin{figure}[H]
\centering
\fbox{\rule{0pt}{0.28\textheight}\rule{0.88\textwidth}{0pt}}
\caption{Research framework}
\label{fig:framework}
\end{figure}

\section{Research design}
\label{sec:design}

\subsection{Sample and data}

The relationship data are from FactSet Revere's Supply Chain Relationships database. We retain CUSTOMER and SUPPLIER records involving a Chinese listed firm and a foreign partner, and preserve the direction of the link. The cleaned source contains 578,142 eligible disclosure records for 2017--2024. Multiple disclosure episodes involving the same Chinese firm, partner, and direction are consolidated into 78,785 distinct relationship pairs. Expanding each pair over the eight calendar years yields the balanced panel of 630,280 relationship--year observations used in the baseline analysis.

Firm financial information is from the China Stock Market and Accounting Research (CSMAR) database. We retain annual consolidated statements, exclude financial firms and ST firms, and winsorize continuous firm variables at the first and 99th percentiles. The matched firm-year file contains 35,978 observations for 5,689 firms (including 3,643 manufacturing firms).

Tariff data combine most-favored-nation rates with official additional-duty schedules from a major non-RCEP import market. Product-level rates are mapped from HS8 products to two-digit national industry categories using a 5,724-row concordance. The industry exposure series uses fixed 2016 Chinese export weights. For the further analysis, firm-level exposure is fixed in 2019 and equals the export-weighted average of the applicable additional duties. Neither measure uses post-policy trade shares.

\subsection{Variable selection and description}

\subsubsection{Supply-chain relationship activity}

The unit of observation is a Chinese firm--foreign partner--direction pair observed at the end of each calendar year. Let $i$ index the Chinese firm, $j$ the foreign partner, and $r$ the direction (supplier or customer). The outcome, $Active_{ijrt}$, equals one if at least one disclosure episode for pair $(i,j,r)$ has started by December 31 of year $t$ and has not ended by that date; it is zero otherwise. Thus, the panel records the extensive margin of relationship activity among pairs observed at least once in the source data, rather than the probability of forming every possible firm combination.

\subsubsection{RCEP treatment}

The treatment is $Post2022_t\times RCEP_j$. In the common-date specification, $Post2022_t$ equals one for calendar years 2022--2024 and zero for 2017--2021. $RCEP_j$ equals one when the partner is located in one of the 14 RCEP economies other than China. This specification treats 2022, when the agreement first entered into force, as the common policy date. A sensitivity specification instead assigns 2023 to Indonesia and the Philippines, whose implementation dates were later. The relationship-pair fixed effect, $\alpha_{ijr}$, absorbs all time-invariant characteristics of a disclosed pair, while year fixed effects absorb common annual shocks.

\subsubsection{Covariates and auxiliary measures}

Table~\ref{tab:summary} reports the variables used in the analysis. The firm covariates are the log of total assets (``Size''), book leverage (``Lev''), return on assets (``ROA''), annual revenue growth (``Growth''), and inventory divided by total assets (``InvDensity''). They are defined at the firm-year level; the preferred fixed-effects estimate is intentionally identified from the treatment interaction and pair and year fixed effects, while these covariates enter the high-dimensional nuisance models in the double-machine-learning sensitivity analysis. The pre-policy relationship-age measure is the number of years from the first disclosed start date to December 31, 2021 for pairs active in 2021. Pre-policy sourcing breadth counts the distinct RCEP partner countries connected to a Chinese firm through supplier links in 2021; the large-firm indicator is based on the median of firms' 2017--2021 mean log assets. Manufacturing status is based on the firm's latest available primary-industry classification observed no later than 2021.

The policy-linked indicators are defined at the partner-country--year level and merged to every relationship involving that country. Incremental tariff relief is the 2017--2019 import-weighted scheduled tariff reduction under the relevant RCEP schedule. The intermediate-goods share is China's imports of goods classified as intermediate inputs under the strict BEC5 definition from partner country $j$, divided by China's total imports from that partner in year $t$. It measures the composition of bilateral imports and, across RCEP members, the degree of regional intermediate-goods trade integration; it is not a measure of firm $i$'s input purchases. The partner-country import share is country $j$'s share of China's total imports in year $t$; a lower value indicates less dependence on that individual source. These aggregates describe policy-relevant conditions and do not measure a firm's preference claim, certificate of origin, transaction value, or overall firm-level sourcing concentration.

\subsection{Model specification}
\label{sec:strategy}

The empirical design exploits the common entry-into-force date while allowing exposure to vary across partner economies. We compare the evolution of relationships involving RCEP partners with that of relationships involving other foreign partners. Pair fixed effects remove time-invariant differences across disclosed firm--partner--direction pairs, and year fixed effects remove shocks common to all relationships. Identification therefore comes from the differential post-2022 change between the two partner groups. The identifying condition is that, absent RCEP, their relationship activity would have followed parallel trends. It also requires the comparison group not to experience differential spillovers from the agreement. That requirement may be imperfect because RCEP can redirect sourcing and sales between member and non-member economies: trade diversion that reduces non-member relationships would enlarge the estimated contrast, whereas positive spillovers to non-members would attenuate it. The coefficient should therefore be interpreted as the post-entry change for RCEP relationships relative to non-RCEP relationships, rather than as an isolated effect against a completely untreated counterfactual. We estimate a linear probability model because the outcome is an indicator and interpret the coefficient as a percentage-point change.

The baseline specification is
\begin{equation}
\label{eq:did}
Active_{ijrt}=\beta(Post2022_t\times RCEP_j)+\alpha_{ijr}+\lambda_t+\varepsilon_{ijrt},
\end{equation}
where $\alpha_{ijr}$ denotes firm--partner--direction pair fixed effects and $\lambda_t$ denotes year fixed effects. The coefficient $\beta$ is the differential change in the probability that an observed relationship is active for RCEP partners relative to other foreign partners. Standard errors are clustered by relationship pair, matching the level at which serially correlated relationship outcomes and treatment histories are observed.

\subsection{Descriptive statistics}

\begin{table}[H]
\centering
\caption{Summary statistics}
\label{tab:summary}
\small
\resizebox{\textwidth}{!}{\begin{tabular}{lrrrrrr}
\toprule
Variable & Observations & Mean & S.D. & Min. & Median & Max. \\
\midrule
\multicolumn{7}{l}{\textit{Relationship--year variables}} \\
Active relationship & 630,280 & 0.226 & 0.418 & -- & -- & -- \\
RCEP partner & 630,280 & 0.183 & 0.386 & -- & -- & -- \\
Post-2022 indicator & 630,280 & 0.375 & 0.484 & -- & -- & -- \\
RCEP $\times$ Post-2022 & 630,280 & 0.068 & 0.253 & -- & -- & -- \\
Foreign supplier relationship & 630,280 & 0.413 & 0.492 & -- & -- & -- \\
Relationship age in 2021 (years) & 630,280 & 0.478 & 1.387 & 0.000 & 0.000 & 18.746 \\
Pre-policy RCEP sourcing breadth & 630,280 & 0.512 & 1.302 & 0.000 & 0.000 & 8.000 \\
Large-firm indicator & 362,440 & 0.501 & 0.500 & -- & -- & -- \\
\midrule
\multicolumn{7}{l}{\textit{Partner-country--year indicators}} \\
Incremental tariff relief (percentage points) & 1,368 & 0.157 & 2.490 & 0.000 & 0.000 & 65.751 \\
Intermediate-goods share of bilateral imports (percent) & 1,355 & 66.811 & 28.800 & 0.000 & 73.156 & 100.000 \\
Partner-country share of China's imports (percent) & 1,376 & 0.579 & 1.469 & 0.000 & 0.035 & 10.882 \\
\midrule
\multicolumn{7}{l}{\textit{Firm-year covariates}} \\
Log total assets & 35,978 & 22.213 & 1.426 & 19.142 & 22.021 & 26.588 \\
Leverage & 35,978 & 0.411 & 0.206 & 0.056 & 0.399 & 0.924 \\
Return on assets & 35,978 & 0.032 & 0.074 & -0.309 & 0.036 & 0.213 \\
Revenue growth & 30,283 & 0.103 & 0.329 & -0.613 & 0.067 & 1.731 \\
Inventory intensity & 35,337 & 0.128 & 0.111 & 0.000 & 0.104 & 0.609 \\
\bottomrule
\end{tabular}
}
\begin{minipage}{0.85\textwidth}
\footnotesize \textit{Notes}: Statistics cover 2017--2024 and use all available observations; counts vary with missing values. Binary variables report the mean and S.D. only, with -- indicating omitted distribution statistics. Percentage measures use the units shown in the variable labels.
\end{minipage}
\end{table}
 
Table~\ref{tab:summary} shows that the balanced relationship--year panel contains 630,280 observations. The smaller counts for the partner-country--year indicators and firm covariates reflect source-file coverage and missing values rather than changes to the baseline sample. Figure~\ref{fig:descriptive} then traces the unconditional active-relationship rates for the two partner groups.

\begin{figure}[H]
\centering
\includegraphics[width=0.88\textwidth]{figures/descriptive_active_links.pdf}
\caption{Descriptive relationship activity by partner group, 2017--2024}
\label{fig:descriptive}
\caption*{\footnotesize Notes: The left plot reports the share of relationship pairs that are active at each year-end, separately for RCEP and other foreign partners. The right plot reports the difference between the two rates in percentage points. The vertical line marks the January 2022 entry into force. The figure is descriptive and does not condition on the fixed effects used in the regressions.}
\end{figure}

Active rates rise for both groups over the sample. The rate for RCEP-partner pairs increases from 12.68 percent in 2017 to 43.14 percent in 2024, while the corresponding rate for other foreign partners rises from 11.72 to 39.68 percent. The raw gap therefore moves from 0.95 percentage points in 2017 to 3.47 percentage points in 2024, with a 1.85-point gap immediately before entry into force. These unconditional patterns motivate the difference-in-differences comparison but are not themselves evidence of a causal effect.

\section{Empirical results and robustness tests}
\label{sec:empirical}

\subsection{RCEP and supply-chain relationship activity}

Table~\ref{tab:baseline} presents the baseline estimates. The preferred coefficient is 0.0136 with a standard error of 0.0040 ($p=0.001$), corresponding to a 1.36 percentage-point differential increase in the active-relationship probability. Using member-specific entry dates gives 0.0144 (standard error 0.0041), a nearly identical estimate. The stability of the point estimate indicates that imposing a common 2022 treatment date is not responsible for the baseline result. The positive and statistically significant estimates are consistent with H1.

\begin{table}[H]
\centering
\caption{Baseline results}
\label{tab:baseline}
\small
\begin{threeparttable}
\resizebox{\textwidth}{!}{\begin{tabular}{lcc}
\toprule
 & (1) Common 2022 date & (2) Member-specific dates \\
 & Pair-clustered S.E. & Pair-clustered S.E. \\
\midrule
RCEP treatment & 0.0136*** & 0.0144*** \\
 & (0.0040) & (0.0041) \\
\midrule
Observations & 630,280 & 630,280 \\
Relationship pairs & 78,785 & 78,785 \\
Pair fixed effects & Yes & Yes \\
Year fixed effects & Yes & Yes \\
\bottomrule
\end{tabular}
}
\begin{tablenotes}[flushleft]
\item
\footnotesize \textit{Notes}: The outcome is an active firm--partner relationship at year-end. All columns include pair and year fixed effects. Column (1) uses a common 2022 entry date; column (2) uses member-specific entry dates. Standard errors are clustered by relationship pair and reported in parentheses. * $p<0.10$, ** $p<0.05$, *** $p<0.01$.
\end{tablenotes}
\end{threeparttable}
\end{table}
 
\subsection{Mechanism analysis}

The theoretical discussion identifies three policy-linked channels associated with relationship activity: incremental tariff relief, regional intermediate-goods trade integration, and reduced dependence on an individual partner-country source. Each measure varies at the partner-country--year level and is assigned to relationship $(i,j,r)$ according to the partner's country. Denoting the relevant country-level measure by $M_{jt}$, we estimate

\begin{align}
M_{jt} &= a(Post2022_t\times RCEP_j)+\alpha_{ijr}+\lambda_t+u_{ijrt}, \label{eq:mechanism_a}\\
Active_{ijrt} &= c'(Post2022_t\times RCEP_j)+bM_{jt}+\alpha_{ijr}+\lambda_t+v_{ijrt}. \label{eq:mechanism_b}
\end{align}

Equations~\eqref{eq:mechanism_a} and \eqref{eq:mechanism_b} define the two paths. For each indicator, the total effect $c$ is re-estimated from Eq.~\eqref{eq:did} on the same complete-case sample used in the two path equations; this is why the reported values of $c$ differ slightly across rows and from the full-sample estimate in Table~\ref{tab:baseline}. The indirect association is $ab$, and the linear decomposition satisfies $c=c'+ab$. Inference for $ab$ is based on 2,000 bootstrap samples drawn at the relationship-pair level; both path equations are re-estimated in every draw. This procedure preserves the panel structure and uses the same fixed effects and clustering level as the corresponding regressions.

Table~\ref{tab:mechanisms} reports the estimates. The percentile 95\% bootstrap confidence intervals exclude zero for all three indirect associations: $[0.0087,0.0209]$ for incremental tariff relief, $[0.0002,0.0004]$ for the intermediate-goods share of bilateral imports, and $[0.0240,0.0307]$ for the partner-country import share. RCEP exposure is associated with a 0.0176 increase in the intermediate-goods share, equivalent to 1.76 percentage points. The corresponding indirect association with relationship activity is 0.0003, or 0.03 percentage points. This estimate indicates stronger intermediate-goods trade integration among RCEP members, not increased input purchases by any particular firm. Combined with the supplier-side estimate reported below, the result is consistent with regional production-network integration helping Chinese firms maintain or establish disclosed supplier relationships. RCEP exposure is also associated with a lower partner-country import share, and a lower share is in turn associated with more active relationships. In this third decomposition, however, the positive indirect association ($ab=0.0274$) is offset by a negative direct coefficient ($c'=-0.0141$); it should not be interpreted as a proportion of the total effect. Overall, the three patterns are consistent with H2 as policy-linked channel evidence, but do not identify firm-level preference use, certification, or transaction-level adjustment.

\begin{table}[H]
\centering
\caption{Policy-linked mechanism evidence}
\label{tab:mechanisms}
\small
\begin{minipage}{\textwidth}
\centering
\smallskip
\resizebox{\textwidth}{!}{\begin{tabular}{lrrrrrr}
\toprule
Mechanism measure & $a$ & $b$ & $c$ & $c'$ & \shortstack{Indirect\\effect ($ab$)} & \shortstack{Bootstrap\\95\% CI} \\
\midrule
Incremental tariff relief & 0.2186*** & 0.0688*** & 0.0133*** & -0.0017 & 0.0150*** & [0.0087, 0.0209] \\
 & (0.0021) & (0.0148) & (0.0040) & (0.0050) & & \\
Intermediate-goods import share & 0.0176*** & 0.0181*** & 0.0151*** & 0.0148*** & 0.0003*** & [0.0002, 0.0004] \\
 & (0.0006) & (0.0035) & (0.0041) & (0.0041) & & \\
Partner-country import share & -0.0127*** & -2.1607*** & 0.0133*** & -0.0141*** & 0.0274*** & [0.0240, 0.0307] \\
 & (0.0001) & (0.1394) & (0.0040) & (0.0043) & & \\
\bottomrule
\end{tabular}
}
\vspace{2pt}
\parbox[t]{\textwidth}{\footnotesize \textit{Notes}: Each row uses its mechanism-specific complete-case sample. The percentile interval for $ab$ is based on 2,000 relationship-pair bootstrap draws. All models include pair and year fixed effects; standard errors in parentheses are clustered by pair. *, **, and *** denote 10\%, 5\%, and 1\% significance, using the corresponding bootstrap intervals for $ab$. All mechanism measures vary at the partner-country--year level.}
\end{minipage}
\end{table}
 
\subsection{Robustness tests}

\subsubsection{Event-study test}

Figure~\ref{fig:event} plots the event-study coefficients relative to the year before entry into force. We test the null hypothesis that the four pre-treatment coefficients for 2017--2020 are jointly equal to zero. The Wald statistic is 8.389 with four degrees of freedom ($p=0.078$). The null is therefore not rejected at the 5\% level but is rejected at the 10\% level. Accordingly, the event study provides only marginal support for the parallel-trends assumption rather than a definitive validation of it.

\begin{figure}[H]
\centering
\includegraphics[width=0.82\textwidth]{figures/parallel_trends_standard_en.pdf}
\caption{Event-study estimates of the RCEP relationship effect}
\label{fig:event}
\caption*{\footnotesize Notes: Points are event-study coefficients from the relationship-pair fixed-effects model with year fixed effects; bars are 95 percent confidence intervals based on relationship-pair clustered standard errors. Event time $-1$ (2021) is the omitted period. The dashed vertical line marks RCEP's entry into force. The joint Wald test evaluates whether all four pre-treatment coefficients are zero.}
\end{figure}

\subsubsection{Alternative treatment timing, samples, and estimator}

Table~\ref{tab:robustness} shows that the pair-clustered estimate remains positive and statistically significant when treatment is assigned using member-specific dates (0.0144), 2020--2021 observations are removed (0.0151), the sample begins in 2019 (0.0107), partners from the largest non-RCEP comparison market are excluded (0.0098), or offshore financial centers are excluded (0.0171). The exclusions reduce reliance on comparison markets that may have distinctive post-2022 dynamics, and their positive estimates indicate that neither group alone drives the baseline result.

We also implement a double machine-learning sensitivity analysis. The outcome is the change in the relationship's active status between the pre- and post-policy periods. Five-fold cross-fitting is grouped by partner country. Gradient-boosting regression and classification models estimate the outcome and treatment nuisance functions using pre-policy relationship history, firm characteristics, and trends. The orthogonalized estimate is 0.0077 (standard error 0.0021, $p=0.002$; 95 percent confidence interval $[0.0034,0.0116]$). Because this procedure uses a pre/post change rather than the annual fixed-effects outcome, it is a functional-form and high-dimensional-control check rather than a mechanically identical re-estimate of $\beta$.

\begin{table}[H]
\centering
\caption{Robustness checks}
\label{tab:robustness}
\scriptsize
\begin{threeparttable}
\begin{tabular*}{\textwidth}{@{\extracolsep{\fill}}lrrr}
\toprule
Specification & Estimate & s.e. & $p$-value \\
\midrule
Member-specific entry dates & 0.0144*** & 0.0041 & $<0.001$ \\
Exclude 2020--2021 observations & 0.0151*** & 0.0047 & 0.001 \\
2019--2024 sample window & 0.0107** & 0.0042 & 0.011 \\
Exclude the largest non-RCEP comparison market & 0.0098** & 0.0042 & 0.019 \\
Exclude offshore financial centers & 0.0171*** & 0.0042 & $<0.001$ \\
Double machine learning & 0.0077*** & 0.0021 & 0.002 \\
\bottomrule
\end{tabular*}
\begin{tablenotes}[flushleft]
\item
\footnotesize \textit{Notes}: Standard errors are clustered by relationship pair. The double-machine-learning row uses five-fold cross-fitting and 999 pair-clustered bootstrap replications. * $p<0.10$, ** $p<0.05$, *** $p<0.01$.
\end{tablenotes}
\end{threeparttable}
\end{table}
 
\section{Heterogeneity analysis}
\label{sec:heterogeneity}

The average estimate can conceal differences in how relationships respond to the agreement. We therefore turn first to the supplier--customer asymmetry predicted by H3 and then examine whether the response varies with pre-policy relationship maturity, prior trade-agreement coverage, and manufacturing exposure. The pooled specifications in Table~\ref{tab:heterogeneity} estimate group-specific treatment coefficients in a common regression and test their equality directly.

\subsection{Relationship direction}

Supplier relationships, which correspond to Chinese firms purchasing from foreign partners, have an estimated effect of 0.0295 ($p<0.001$), while customer relationships have an estimate of -0.0033 ($p=0.531$). The difference is 0.0327 and is statistically significant at the 1\% level. The concentration of the response in supplier links is consistent with regional cumulation operating most directly through intermediate-input procurement and supports H3.

\subsection{Pre-policy relationship maturity}

The scope for policy-induced adjustment may also depend on relationship-specific capital accumulated before implementation. Long-standing links have already absorbed search, screening, negotiation, and adaptation costs, leaving less room for an institutional change to alter their annual activity. We therefore classify a relationship as long-standing when it was active at year-end 2021 and had accumulated at least three years of relationship age by that date. The three-year threshold corresponds to the complete 2019--2021 pre-policy window and is fixed independently of the estimated coefficients. The remaining disclosed pairs form the comparison group. Both group-specific effects are estimated jointly in the baseline specification. The estimate for other relationships is 0.0183 ($p<0.001$), whereas the estimate for long-standing relationships is -0.0162 and statistically insignificant ($p=0.309$). The difference is -0.0345 ($p=0.036$), indicating that the positive response is concentrated among relationships with greater scope for adjustment rather than among mature incumbent links.

\subsection{Prior trade-agreement coverage}

RCEP did not provide the same degree of institutional novelty across members. China already had free-trade agreements with the ASEAN members, Australia, New Zealand, and Korea, whereas RCEP created the first free-trade agreement between China and Japan. We therefore compare Japan with the RCEP partners covered by an earlier agreement, using the same set of non-RCEP relationships as the common comparison group. The coefficient for partners with prior agreement coverage is -0.0026 ($p=0.605$), while the coefficient for Japan is 0.0378 ($p<0.001$). Their difference is 0.0403 ($p<0.001$). This pattern is consistent with a larger initial response where RCEP introduced genuinely new bilateral market-access rules, although the Japan estimate also reflects characteristics specific to a single partner economy.

\subsection{Manufacturing exposure among supplier links}

Regional cumulation and origin requirements are most directly relevant to firms that organize manufacturing inputs across borders. We classify firms using the latest primary-industry code observed no later than 2021 and restrict this comparison to supplier relationships. The estimated effect is 0.0274 for manufacturing firms ($p=0.030$), compared with -0.0205 for non-manufacturing firms ($p=0.263$). The difference is 0.0479 ($p=0.031$). The supplier-side response is therefore concentrated among manufacturing firms, reinforcing the input-sourcing rationale for the directional asymmetry in H3.

\begin{table}[H]
\centering
\caption{Heterogeneity analysis}
\label{tab:heterogeneity}
\small
\begin{minipage}{\textwidth}
\centering
\resizebox{\textwidth}{!}{\begin{tabular}{lccc}
\toprule
Partition & Group 1 & Group 2 & Difference (Group 2 $-$ Group 1) \\
\midrule
Relationship direction & Customer: -0.0033 & Supplier: 0.0295*** & 0.0327*** \\
 & (0.0052) & (0.0063) & (0.0082) \\
Pre-policy relationship maturity & Other: 0.0183*** & Long-standing: -0.0162 & -0.0345** \\
 & (0.0041) & (0.0160) & (0.0165) \\
Pre-RCEP FTA coverage & Existing FTA: -0.0026 & Japan: 0.0378*** & 0.0403*** \\
 & (0.0050) & (0.0061) & (0.0075) \\
Firm sector, supplier links & Non-manufacturing: -0.0205 & Manufacturing: 0.0274** & 0.0479** \\
 & (0.0183) & (0.0127) & (0.0223) \\
\bottomrule
\end{tabular}
}
\vspace{2pt}
\parbox{\textwidth}{\footnotesize \textit{Notes}: Group coefficients are estimated jointly. A long-standing relationship is active at year-end 2021 and is at least three years old. The sector comparison is restricted to supplier links with pre-policy industry data. All models include relationship-pair and year fixed effects; relationship-pair clustered standard errors are in parentheses. The difference is Group 2 minus Group 1. * $p<0.10$, ** $p<0.05$, *** $p<0.01$.}
\end{minipage}
\end{table}
 
\section{Further analysis}
\label{sec:further}

\subsection{Pre-existing external tariff exposure}

The robustness analysis shows that the estimated contrast is sensitive to the composition of the non-RCEP comparison group, a concern related to spillovers and trade diversion generated by RCEP itself. A separate alternative explanation is that the post-2022 pattern reflects firms' reorientation following pre-existing tariff pressure outside the agreement. To assess this possibility, we interact RCEP exposure with a firm-level measure of additional-duty exposure fixed in 2019 and estimate Eq.~\eqref{eq:triple}:
\begin{equation}
\label{eq:triple}
\begin{aligned}
Active_{ijrt}={}&\beta_1(Post2022_t\times RCEP_j)+\beta_2(Post2022_t\times TariffExposure_i)\\
&+\beta_3(Post2022_t\times RCEP_j\times TariffExposure_i)\\
&+\alpha_{ijr}+\lambda_t+\varepsilon_{ijrt}.
\end{aligned}
\end{equation}
The product weights underlying the exposure measure are fixed in 2016, and the firm-level measure is fixed before RCEP entered into force. A positive $\beta_3$ would indicate stronger reorientation toward RCEP partners among firms with greater pre-existing exposure.

Table~\ref{tab:tariffalternative} reports estimates on the subsample for which 2019 firm-level tariff exposure can be matched. The first row is therefore a matched-sample baseline, rather than a replication of the full-sample estimate in Table~\ref{tab:baseline}; the restriction leaves 253,184 observations and 31,648 relationship pairs. The matched-sample coefficient is 0.0061 (standard error 0.0065, $p=0.346$) and changes little after adding the post-period tariff-exposure control (0.0063, standard error 0.0065, $p=0.331$). The triple interaction is 0.0471 (standard error 0.0606, $p=0.437$). Thus, firms with higher pre-existing tariff exposure do not exhibit a statistically larger differential change in RCEP relationships after 2022. The evidence does not support this external-tariff reorientation explanation, but the matched sample covers only 40.2 percent of relationship pairs and the imprecision of the estimates prevents a strict exclusion claim.

\begin{table}[H]
\centering
\caption{Further analysis: Pre-existing external tariff exposure}
\label{tab:tariffalternative}
\small
\begin{threeparttable}
\begin{tabular}{lrrr}
\toprule
Term & Estimate & s.e. & $p$-value \\
\midrule
Post2022 $\times$ RCEP (matched baseline) & 0.0061 & 0.0065 & 0.346 \\
Post2022 $\times$ RCEP (matched + tariff control) & 0.0063 & 0.0065 & 0.331 \\
Post2022 $\times$ RCEP $\times$ tariff exposure & 0.0471 & 0.0606 & 0.437 \\
\bottomrule
\end{tabular}
\begin{tablenotes}[flushleft]
\item
\footnotesize \textit{Notes}: The sample is restricted to firm--partner relationships for which the 2019 firm-level tariff exposure can be matched (253,184 observations and 31,648 relationship pairs). All models include relationship-pair and year fixed effects, with standard errors clustered by relationship pair. The triple interaction tests whether high-exposure firms experienced an additional change in relationships with RCEP partners.
\end{tablenotes}
\end{threeparttable}
\end{table}
 
\section{Conclusion and policy recommendations}
\label{sec:conclusion}

This paper examines how RCEP's entry into force relates to the international supply-chain relationships of Chinese listed firms. In a panel of 630,280 relationship--year observations and 78,785 relationship pairs, the preferred specification yields a 0.0136 differential increase in the probability that an RCEP relationship is active at year-end. The corresponding estimate is 0.0144 under member-specific implementation dates. The response is concentrated in supplier relationships, whose estimated increase is 0.0295, compared with an insignificant estimate of -0.0033 for customer relationships. It is also larger for Japan than for members with prior free-trade-agreement coverage (a difference of 0.0403) and, among supplier links, for manufacturing than for non-manufacturing firms (a difference of 0.0479). At the partner-country--year level, the bootstrap indirect associations are 0.0150 for incremental tariff relief (95\% CI $[0.0087,0.0209]$), 0.0003 for the intermediate-goods share of bilateral imports (95\% CI $[0.0002,0.0004]$), and 0.0274 through a lower partner-country import share (95\% CI $[0.0240,0.0307]$). These aggregate patterns, together with the supplier-side response, are consistent with RCEP strengthening regional production-network integration and thereby supporting the maintenance or formation of disclosed supplier relationships.

The findings carry three policy implications. First, implementation efforts should give firms clear and accessible guidance on regional cumulation, tariff schedules, and documentary requirements. The value of a common agreement depends not only on formal concessions but also on whether firms can apply the rules at reasonable administrative cost. Second, the increase in the member-level intermediate-goods share, together with the supplier-side relationship response, suggests that trade-facilitation policies should support the matching of firms with qualified regional suppliers. Third, policymakers should monitor whether the early expansion of regional links reduces dependence on individual sources while preserving resilience to disruptions. These measures would help translate formal integration into durable firm-level relationships.

The evidence should nevertheless be interpreted within the limits of the design. The joint Wald test of the four pre-treatment event-study coefficients does not reject at the 5\% level but is marginal at the 10\% level ($\chi^2(4)=8.389$, $p=0.078$). The non-RCEP comparison group may itself respond to trade diversion or production-network spillovers generated by RCEP and is therefore not guaranteed to provide a completely untreated counterfactual. The estimates remain positive after excluding the largest non-RCEP comparison market and offshore financial centers. FactSet Revere primarily covers listed and large firms and does not report transaction values, customs declarations, certificates of origin, or firms' actual preference utilization. Moreover, the 2022--2024 period captures only the initial phase of a tariff schedule extending over several years. Linking product-level customs records to firm identifiers and extending the post-entry window would allow future research to distinguish relationship formation, survival, and trade-value responses more precisely.

\bibliographystyle{elsarticle-harv}
\begin{thebibliography}{99}

\bibitem[Ali et al.(2026)]{ali2026governance}
Ali, W., Saini, S., Dash, D.P., 2026. Institutional governance and energy security: Unveiling the drivers of global value chain participation in Asian economies. Journal of Asian Economics 103, 102108.

\bibitem[Baier and Bergstrand(2007)]{baier2007free}
Baier, S.L., Bergstrand, J.H., 2007. Do free trade agreements actually increase members' international trade? Journal of International Economics 71(1), 72--95.

\bibitem[Crosignani et al.(2022)]{crosignani2022pirates}
Crosignani, M., Macchiavelli, M., Silva, A.F., 2022. Pirates without borders: The propagation of cyberattacks through firms' supply chains. Journal of Financial Economics 147(2), 432--456.

\bibitem[Du et al.(2023)]{du2023higher}
Du, Y., Yan, J., Cao, F., Li, Y., Zhou, M., 2023. Higher education expansion and domestic value added in exports: Theory and evidence from China. Journal of Asian Economics 87, 101621.

\bibitem[Garg et al.(2023)]{garg2023within}
Garg, S., Ghosh, P., Tan, B.J., 2023. Within firm supply chains: Evidence from India. Journal of International Economics 144, 103793.

\bibitem[Grossman et al.(2024)]{grossman2024tariffs}
Grossman, G.M., Helpman, E., Redding, S.J., 2024. When tariffs disrupt global supply chains. American Economic Review 114(4), 988--1029.

\bibitem[Handley et al.(2024)]{handley2024input}
Handley, K., Lim\~ao, N., Ludema, R.D., Yu, Z., 2024. Firm input choice under trade policy uncertainty. Journal of International Economics 150, 103909.

\bibitem[Hayakawa(2022)]{hayakawa2022multiple}
Hayakawa, K., 2022. Multiple preference regimes and rules of origin. Review of World Economics 159(3), 673--696.

\bibitem[Hayakawa et al.(2023)]{hayakawa2023firm}
Hayakawa, K., Laksanapanyakul, N., Yoshimi, T., 2023. Firm-level utilization rates of regional trade agreements: Importers' perspective. Journal of Asian Economics 86, 101610.

\bibitem[Hayakawa and Yamanouchi(2023)]{hayakawa2023reform}
Hayakawa, K., Yamanouchi, K., 2023. How does the reform of rules of origin affect firm performance in importing countries? The World Economy 47(2), 837--856.

\bibitem[Hayakawa et al.(2024)]{hayakawa2024negotiating}
Hayakawa, K., Kimura, F., Mukunoki, H., Urata, S., 2024. Negotiating over the rules of origin in regional trade agreements in Asia. The World Economy 47(7), 3097--3119.

\bibitem[He et al.(2023)]{he2023regional}
He, Y., Huo, W., Yu, J., 2023. Tracing the regional dual value chains: Measurement on the production position and evidence from China. Journal of Asian Economics 85, 101591.

\bibitem[Head et al.(2024)]{head2024laffer}
Head, K., Mayer, T., Melitz, M., 2024. The Laffer curve for rules of origin. Journal of International Economics 150, 103911.

\bibitem[Huang et al.(2026)]{huang2026intermediate}
Huang, F., Li, J., Shen, J., Wang, H., 2026. Digital economy development and repositioning in the global intermediate goods trade network. Journal of Asian Economics 105, 102218.

\bibitem[Kiyota(2022)]{kiyota2022covid}
Kiyota, K., 2022. The COVID-19 pandemic and the world trade network. Journal of Asian Economics 78, 101419.

\bibitem[Lee et al.(2023)]{lee2023thirdcountry}
Lee, W., Mulabdic, A., Ruta, M., 2023. Third-country effects of regional trade agreements: A firm-level analysis. Journal of International Economics 140, 103688.

\bibitem[Li et al.(2025)]{li2025integration}
Li, Y., Wan, Y., Xie, S., 2025. The employment effects of regional market integration: Evidence from China. Journal of Asian Economics 101, 102053.

\bibitem[Li and Li(2026)]{li2026digitalrules}
Li, N., Li, R., 2026. A study on the impact of regional trade agreement digital trade rules on global innovation network. Journal of Asian Economics 103, 102144.

\bibitem[Liu et al.(2026)]{liu2026risks}
Liu, Y., Deng, X., Min, J.H., Wang, J., Liu, Y., 2026. Multidimensional risks and the cleaner production-oriented trade transformation in the China--ASEAN. Journal of Asian Economics 103, 102128.

\bibitem[Liu and Zhen(2022)]{liu2022rcep}
Liu, B., Zhen, Y., 2022. Trade and rule effects of RCEP. World Economy Studies.

\bibitem[Ma et al.(2024a)]{ma2024data}
Ma, S., Huang, S., Wu, P., 2024a. Data policy restrictions and cross-border e-commerce: Evidence from China. Journal of Asian Economics 95, 101826.

\bibitem[Ma et al.(2024b)]{ma2024mobility}
Ma, S., Ren, F., Feng, F., Feng, K., 2024b. Cross-border personnel mobility and bilateral value chain linkages: Evidence from visa liberalization in China. Journal of Asian Economics 95, 101844.

\bibitem[Narayan et al.(2022)]{narayan2022agreements}
Narayan, P.K., Rath, B.N., Syarifuddin, F., 2022. Understanding the role of trade agreements in Indonesia's FDI. Journal of Asian Economics 82, 101532.

\bibitem[Newman et al.(2023)]{newman2023imports}
Newman, C., Rand, J., Tarp, F., 2023. Imports, supply chains and firm productivity. World Development 172, 106371.

\bibitem[Ni et al.(2022)]{ni2022standards}
Ni, S., Bai, X., Li, L., 2022. Behind the heterogeneous trade effects of standards: Multi-sector evidence from China. Journal of Asian Economics 82, 101513.

\bibitem[Ou et al.(2026)]{ou2026spillovers}
Ou, J., Yang, T., Yao, S., Hu, M., 2026. Digital finance, supply chain spillovers and corporate innovation: Evidence from Chinese listed firms 2011--22. Journal of Asian Economics 104, 102175.

\bibitem[Peng et al.(2025)]{peng2025financing}
Peng, Q., Chen, Z., Li, J., 2025. Financing constraints and the export partnerships with RCEP: Evidence from manufacturing firms in China. Journal of Asian Economics 97, 101870.

\bibitem[Ping(2021)]{ping2021rcep}
Ping, L., 2021. Institutional characteristics and economic effects of RCEP. Asia-Pacific Economic Review.

\bibitem[Poilly and Tripier(2025)]{poilly2025regional}
Poilly, C., Tripier, F., 2025. Regional trade policy uncertainty. Journal of International Economics 155, 104078.

\bibitem[Shen and Li(2020)]{shen2020rcep}
Shen, M., Li, T., 2020. RCEP: Institutional framework, rules, and regional effects. International Economic Review.

\bibitem[Sun et al.(2022)]{sun2022rtaquality}
Sun, J., Luo, Y., Zhou, Y., 2022. The impact of regional trade agreements on the quality of export products in China's manufacturing industry. Journal of Asian Economics 80, 101456.

\bibitem[Sun et al.(2023)]{sun2023estimating}
Sun, K., Xiao, H., Jia, Z., Tang, B., 2023. Estimating the effects of regional value chains of the RCEP in a GVC-CGE model. Journal of Asian Economics 88, 101647.

\bibitem[Tanaka and Fukunishi(2022)]{tanaka2022rooexports}
Tanaka, K., Fukunishi, T., 2022. Rules of origin and exports in developing economies: The case of garment products. Journal of Asian Economics 82, 101514.

\bibitem[Wu et al.(2023)]{wu2023gvc}
Wu, J., Bian, Y., Bai, J., 2023. Global value chain upgrading effect of foreign capital withdrawal: Evidence from Chinese firms. Journal of Asian Economics 89, 101672.

\bibitem[Xie(2023)]{xie2023tariff}
Xie, D., 2023. Tariff cost and cross-border M\&A affiliate sales: Evidence from China. Journal of Asian Economics 87, 101636.

\bibitem[Xu et al.(2024)]{xu2024financial}
Xu, L., Jing, J., Wu, C., 2024. Impact of financial development on the position in global value chain: An analysis from the perspective of R\&D intensity. Journal of Asian Economics 92, 101742.

\bibitem[Zhang and Ding(2024)]{zhang2024restructuring}
Zhang, H., Ding, Y., 2024. The impact of global value chain restructuring on the product structure of China--EU exports. Journal of Asian Economics 91, 101708.

\bibitem[Zhou and Wen(2022)]{zhou2022trade}
Zhou, F., Wen, H., 2022. Trade policy uncertainty, development strategy, and export behavior: Evidence from listed industrial companies in China. Journal of Asian Economics 82, 101528.

\bibitem[Zhou and Wu(2025)]{zhou2025origin}
Zhou, H., Wu, L., 2025. Rules of origin, import of inputs and firm innovation. Journal of Asian Economics 96, 101872.

\bibitem[Zhu et al.(2025)]{zhu2025diversifying}
Zhu, J., Liu, X., Cao, L., 2025. The Belt and Road Initiative and China's diversifying food imports: A perspective on the ``Five-pronged approach.'' Journal of Asian Economics 96, 101849.

\bibitem[Zhu and Sun(2025)]{zhu2025environmental}
Zhu, Y., Sun, C., 2025. Environmental provisions in preferential trade agreements and export product quality of Chinese firms. Journal of Asian Economics 101, 102048.

\end{thebibliography}

\end{document}
